\section{A replacement theorem with feedback}
\label{feedback}
A replacement can have writable, unbounded data and still admit a uniform decision procedure for every input. The following result makes that assertion precise. It also specifies the operation unavailable to this class: a polynomial update can copy an entire old register, but cannot expose one of its arbitrary interior subtrees as a new register.

Let \(Q\) be any finite algebra for a constant \(s\) and a binary operation \(\mu\). Write \(h:\T\to Q\) for its unique homomorphism. Taking products allows any fixed finite family of regular tree observations simultaneously. Choose \(r\) tree registers, an arbitrary function \(d:Q^r\to \T(x_1,\dots,x_r)^r\) and \(A\subseteq Q^r\). At a tick, evaluate the expressions \(d(h^r(\mathbf t))\) simultaneously in the old registers. Variables may occur any finite number of times. The register values can grow exponentially in the number of ticks even from a single small seed. Each fixed machine has a bounded per-tick growth factor. Acceptance reads \(A\) after any number of ticks, including zero. Initial registers may be computed by any total effective encoder.

\begin{theorem}\label{cutoff}
Write \(q=|Q|\). A replacement reaches \(A\) if and only if it does so at some \(k<q^r\). The induced halting domain is decidable. With a regular-preimage encoder, the halting domain on the original inputs is regular.
\end{theorem}
\begin{proof}
For each expression \(e\), structural induction gives
\[
 h(e_{\T}(\mathbf t))=e_Q(h^r(\mathbf t)).
\]
For a variable this is projection, for a constant it is the value \(s\), and for an application the induction hypotheses compose with \(\mu\). Thus the vector observation after one concrete tick is
\(\delta(a)=d(a)_Q(a)\), where \(a=h^r(\mathbf t)\). Induction gives equality after every number of ticks, even though the choice of expression depends on every preceding observation. If an earliest accepting index were at least \(q^r\), its prefix including that index would have two equal abstract states. Removing the segment between these states preserves the suffix up to acceptance, contradicting minimality. Enumerate the bounded orbit to decide. For the language claim, the acceptable initial vectors form a finite subset of \(Q^r\), so pull that subset back along the encoder.
\end{proof}

The bound concerns a joint observation state for an individual register. A finite external controller of size \(m\) can be included by taking a product on the vector orbit, giving \(m q^r\). A decoder may additionally count designated output occurrences after acceptance. As long as acceptance alone is equivalent to source halting, an unbounded numerical output cannot invalidate this decision procedure. The theorem does not include arbitrary rewriting at an observed address. In our construction a slot is opened and the designated continuation of its new word becomes the subject of later updates. On the one-word program \((1)\) with seed \(1\), every pop after the first is at the residual exposed by the immediately preceding firing. Induction on these firings proves this for all times. The position of the active residual descends through frames of unbounded depth. A clock presenting candidate numbers to an arithmetic history verifier has no such successor invariant. Applying the cutoff to such a clock requires a finite joint observation algebra for its verifier. The old arithmetic acceptance test has no such representation on the clock inputs that would give a decidable universal domain.

\begin{prop}\label{stack}
A finite controller on one unbounded pushdown stack, with regular whole-stack guards and a regular halting test, has decidable existential halting from every effectively constructed finite initial configuration. This also holds if guarded moves are nondeterministic.
\end{prop}
\begin{proof}
Use a product DFA for the guards and the target test, reading from the bottom. Decorate each stack symbol with the DFA state for the suffix below it. The state for the whole stack is obtained from the top and its decoration. A pop exposes the previous summary. A push installs the current summary below its new symbol. An endmarker handles emptiness. Finite top replacements maintain valid decorations by finitely many moves, giving an ordinary pushdown system with a regular target.

For completeness, form a finite automaton for that target with the pushdown controls as allowed starting states. Copy its internal automaton into disjoint states and take no initial edges entering controls. For each rule \(pA\to q u\) and existing path \(q\xrightarrow{u}v\), insert \(p\xrightarrow{A}v\). This uses finitely many possible edges, hence terminates. At the fixed point each actual predecessor of a recognized configuration is recognized: decompose its successor's accepting path after the prefix \(u\). Induction on execution length gives completeness. For soundness, in a saturated accepting path from a control, a first edge from the original automaton starts a wholly original accepting path since no added edge leaves an internal state. Otherwise replace its first, added edge by the earlier edges that justified it, and take the associated pushdown move. This strictly decreases the multiset of addition ranks along the accepting path, even if the replacement path is longer. Eventually the remaining accepting path is original, including the possible empty path. Thus the saturated automaton recognizes exactly the predecessor set. Test the initial decorated stack in it \cite{pushdown}.
\end{proof}

The two replacement theorems grant different operations. The first permits copying whole trees; the second allows popping and inspecting nested symbols. Neither covers a selector coupled to arbitrary occurrence-specific copying rewrites of a binary tree. These bounds exclude clock verification and both register and stack replacements with the stated observations. In the selected reduction, the walker carries the consumed bit to the next slot. Within this layout, moving that routing entirely into a fixed sequence confined to the slot is impossible when the stored word is nonempty. The two popped-bit cases have identical unused slots. Any identical sequence of addressed operations internal to those slots produces identical released payloads, whereas the two required outputs are \(\varepsilon\) and \(v\). A construction that joins the consumed cell and the slot in a new redex would fall outside this locality argument.

\section{Size, navigation and extraction}
Let \(N=2L-1\) denote unfolded nodes for a tree with \(L\) leaves; let \(h\) be its edge height. The direct encoder prints a constant number of program templates and one initial word. For input presented as a list of binary words, it takes \(O(1+p+|w|+\sum_i |v_i|)\) output nodes and time. Constructors add a bounded number of nodes per slot and letter. A finite-copying MSO interpretation suffices on a delimited input word: give each letter and delimiter a fixed finite set of copies for its local constructor, link consecutive cell and slot roots within their respective list, and attach the distinguishable endpoints to the constant seed nodes. Endpoints and next delimiters are first-order tests on the ordered input; fixed copies provide each occurrence of the program. This is an output-construction argument, not an invocation of the source interpreter. For a fixed numerical template the bound is \(O(1+e+x)\) with unary inputs. It is exponential in the worst case if \(e,x\) are supplied in binary. Printing the constant numerical menu has a substantial one-time cost. No such linear construction bound is claimed for converting an arbitrary Turing-machine description to the upstream fixed Cook input.

With parent and incoming-side navigation on the unfolded occurrence tree, the walker has finite control, one read-only cursor and no writable stack. A machine storing the cursor as an index into an explicit table needs \(O(\log(N+1))\) bits for that index; parent links and incoming sides belong to the input representation. With downward-only immutable references the executable implements a zipper using \(O(h)\) references for the current context. Storing a path as a repeatedly concatenated string can cost quadratic time in its length. A strict finite tree automaton for the detector uses constant information per node and a postorder implementation needs \(O(h)\) stack frames. The numerical decoder emits one symbol per designated live unit using the parsed roles; a binary count needs \(O(\log(N+1))\) result bits. Even if implemented by an explicit stack or temporary annotations, a scan reads only the current finite term. The operations of the model are independent of such implementation choices.

For a contraction with third argument \(z\), \(L'=L-1+L(z)\). The local height increases by at most one. These facts persist in a surrounding context (take sums for leaves and maxima for height). Induction gives \(L_k\le 2^k L_0\) and \(h_k\le h_0+k\). Hence repeated linear tree scans have total cost \(O((2^k-1)N_0+k)\) through \(k\) contractions, apart from initialization. Fresh copies of the full term need \(O(2^k N_0)\) peak tree space if old trees can be released. Immutable DAG evaluation may reuse the two references to \(z\). For an address of depth \(d\), allocate at most three application nodes for the right side and \(d\) copied ancestors. Unfolding this rebuild of the unique occurrence gives the exact contextual tree replacement, by induction on the address. Even retaining every previous root, the number of allocated DAG nodes is at most
\[
 D_0+3k+k h_0+k(k-1)/2 .
\]
This is a storage bound for nodes constructed by contraction and path copying. It excludes selector workspaces and memo tables, and it gives no corresponding polynomial runtime for scanning the unfolded tree. The two-to-four count excludes navigation, printing the seed, conversion from another source model, the number of letters per register sweep, and cost of materializing an output.
